\chapter{Bowen's Disease}

\section*{Definition and Overview}
Bowen's disease is an early form of skin cancer confined entirely to the top layer of the skin, the epidermis. It is also called squamous cell carcinoma in situ, meaning that the abnormal keratinocytes have not yet invaded the deeper layers of the skin. Bowen's disease appears as a slow-growing, scaly, red-brown patch or plaque that may be mistaken for eczema, psoriasis, or a fungal infection. It is not dangerous in itself, does not spread to lymph nodes or other organs, and is highly treatable. However, without treatment, around 3--5\% of patches progress over time to an invasive squamous cell carcinoma, which is capable of spreading. A related condition affecting the moist skin of the vulva or penis is called erythroplasia of Queyrat.

\section*{Epidemiology}
Bowen's disease most often affects fair-skinned people over the age of 60, and it is slightly more common in women than in men, in part because the lower legs are a particularly common site in women. It develops on sun-exposed skin, most often the lower legs, head and neck, and arms, but it can occur anywhere, including on areas that have had limited sun exposure. People with a history of heavy sun exposure, solarium use, immunosuppressive medicines (such as after organ transplantation), or previous skin cancers are at higher risk. With ageing populations and cumulative sun exposure over decades, the incidence is rising, and many affected people develop more than one patch over their lifetime.

\section*{Aetiology and Pathophysiology}
\begin{itemize}
    \item \textbf{Ultraviolet radiation:} cumulative damage from sun exposure is the dominant risk factor; UV light causes mutations in skin cells that drive abnormal growth.
    \item \textbf{Solariums:} artificial UV exposure substantially increases the risk of Bowen's disease and other skin cancers.
    \item \textbf{Immunosuppression:} medicines taken after organ transplantation or for autoimmune disease impair the skin's ability to repair UV damage, increasing both risk and recurrence.
    \item \textbf{Human papillomavirus (HPV):} certain HPV types have been associated with some cases, particularly in the anogenital region.
    \item \textbf{Arsenic exposure:} rarely, past exposure to contaminated water or medicines contributes to the disease.
    \item \textbf{Mechanism:} the abnormal cells are confined to the epidermis above the basement membrane, so they do not reach the blood vessels and lymphatics beneath; this containment is what makes the disease non-invasive, until a small proportion of lesions breach the membrane and become invasive.
    \item \textbf{Chronic sun damage:} the surrounding skin typically shows other signs of photodamage, such as actinic keratoses and pigmentation, reflecting the field of damage in which Bowen's disease develops.
\end{itemize}

\section*{Clinical Features}
A patch of Bowen's disease typically:
\begin{itemize}
    \item Grows slowly over months to years
    \item Is a flat or slightly raised, scaly, red-brown plaque with well-defined but irregular edges
    \item Ranges in size from a few millimetres to several centimetres
    \item May look like eczema, psoriasis, or a fungal infection that fails to clear with treatment
    \item Can crust, fissure, or bleed if scratched
    \item Is usually painless and often causes no symptoms at all
    \item Occurs most often on the lower legs, but can appear on any skin surface
    \item In the anogenital form, presents as a red, velvety patch that may be itchy or painful
\end{itemize}

\section*{Differential Diagnosis}
Many other skin conditions resemble Bowen's disease:
\begin{itemize}
    \item \textbf{Eczema and psoriasis:} scaly, erythematous patches that respond to topical steroid treatment.
    \item \textbf{Superficial basal cell carcinoma:} a slow-growing, red scaly lesion with a thread-like pearly border.
    \item \textbf{Actinic keratosis:} smaller, rougher "sun spots" that are less persistent and also have malignant potential.
    \item \textbf{Seborrhoeic keratosis or viral warts:} benign growths with a waxy or verrucous surface.
    \item \textbf{Fungal skin infection (tinea):} an annular scaly rash that can closely mimic Bowen's disease.
    \item \textbf{Pigmented Bowen's disease:} a rare variant that can resemble melanoma or an atypical mole.
    \item \textbf{Invasive squamous cell carcinoma:} a raised, indurated, or ulcerating lesion in which invasion must be excluded.
\end{itemize}

\section*{When to Seek Medical Care}
See a doctor if you notice any new, scaly, or crusting skin patch that does not heal, slowly enlarges, or bleeds, especially if you are fair-skinned or have had significant sun exposure. Any skin lesion that is new, changing, or different from the surrounding skin deserves a professional check.

\textbf{Seek urgent medical attention for:}
\begin{itemize}
    \item A known patch that suddenly grows, becomes raised, or develops a lump
    \item A patch that ulcerates, bleeds easily, or becomes painful
    \item Spreading redness, swelling, or discharge, which may indicate infection
    \item A rapidly changing skin lesion of any kind
\end{itemize}

\section*{Diagnosis and Investigations}
\begin{itemize}
    \item \textbf{Visual examination:} a doctor inspects the lesion, usually with a magnifying lens or dermatoscope, and assesses the surrounding skin.
    \item \textbf{Skin biopsy:} a small sample taken under local anaesthetic and examined under a microscope confirms the diagnosis and, importantly, excludes invasive cancer.
    \item \textbf{Punch or shave biopsy:} the standard sampling methods; a shave biopsy can remove many smaller lesions entirely.
    \item \textbf{Full skin check:} examination of the whole skin surface to look for other sun-damaged areas, actinic keratoses, and possible skin cancers elsewhere.
    \item \textbf{History:} past sun exposure, solarium use, immunosuppressant medicines, and previous skin cancers all inform the diagnosis and management.
\end{itemize}

\section*{Management}
Treatment is usually straightforward and depends on the site, size, number, and suspicion of invasion:
\begin{itemize}
    \item \textbf{Curettage and cautery:} scraping away the abnormal tissue and sealing the base with heat; a common, effective first choice, particularly for the lower legs.
    \item \textbf{Cryotherapy:} freezing the patch with liquid nitrogen, often used for small lesions.
    \item \textbf{Topical creams:} imiquimod or fluorouracil cream applied for several weeks, useful for large, multiple, or awkwardly placed patches.
    \item \textbf{Photodynamic therapy:} a light-sensitising cream activated by a special light source, used for the face and scalp where cosmetic results matter.
    \item \textbf{Surgical excision:} cutting out the lesion with a margin of normal skin, preferred when invasive cancer cannot be excluded.
    \item \textbf{Radiotherapy:} occasionally used for lesions that are difficult to treat otherwise, particularly in older people.
    \item \textbf{Follow-up:} review to confirm healing, and long-term surveillance, because new patches and other skin cancers are common in affected people.
\end{itemize}

\section*{Complications}
Complications are uncommon but possible:
\begin{itemize}
    \item Progression to invasive squamous cell carcinoma in around 3--5\% of untreated patches over time
    \item Recurrence of the patch after treatment, more likely with non-surgical methods
    \item Scarring, changes in skin colour, or delayed healing at the treatment site
    \item The appearance of new patches of Bowen's disease elsewhere on the skin
    \item An increased overall risk of other sun-related skin cancers, including melanoma
    \item In the anogenital form, a somewhat higher rate of progression to invasive disease
\end{itemize}

\section*{Prognosis and Outcomes}
The outlook is excellent. Bowen's disease is confined to the top layer of skin, does not spread to other organs, and treatment is very effective, with most lesions cured by a single course of treatment and recurrence rates low. Even when a patch does progress to invasive squamous cell carcinoma, it is usually detected and treated early, with an excellent outcome. The main ongoing concern is not the original patch itself but the increased risk of further sun-related skin lesions, so regular skin checks, sun protection, and prompt attention to any new or changing lesion are the most important long-term measures.

\section*{Prevention and Self-Care}
\begin{itemize}
    \item Protect your skin from the sun: wear protective clothing and a broad-brimmed hat, and apply SPF50+ sunscreen.
    \item Avoid sun exposure in the middle of the day when UV is strongest, and seek shade outdoors.
    \item Never use solariums.
    \item Check your skin regularly for new or changing patches, and ask a partner or doctor to check hard-to-see areas such as the back and scalp.
    \item See a doctor promptly about any scaly patch that does not clear with moisturiser or simple treatment.
    \item Attend recommended skin checks, especially if you are fair-skinned, immunosuppressed, or have had skin cancers before.
    \item Follow the follow-up plan after treatment so that recurrences are caught early.
\end{itemize}

\section*{Sources and Further Reading}
The following sources provide reliable, current information on Bowen's disease and related skin conditions.
\begin{itemize}
    \item NHS. \textit{Bowen's disease.} https://www.nhs.uk/conditions/bowens-disease/
    \item Mayo Clinic. \textit{Squamous cell carcinoma of the skin.} https://www.mayoclinic.org/
    \item Centers for Disease Control and Prevention. \textit{Skin cancer prevention and information.} https://www.cdc.gov/cancer/skin/
    \item World Health Organization. \textit{Ultraviolet radiation and skin cancer prevention.} https://www.who.int/
    \item NICE. \textit{Skin cancer: recognition and referral.} https://www.nice.org.uk/
    \item MedlinePlus. \textit{Skin cancer.} https://medlineplus.gov/
\end{itemize}
